\chapter{Tasks in the Company}
The internship includes a broad range of side tasks in addition to the main internship project, providing real-world working experience and insight into office culture.

\section{Meetings}

\subsection{OPSDEV Daily Stand-up}
Every morning at 09:05, a short meeting is held via Google Meet where team members share updates on their current work and discuss ongoing tasks.
\subsection{OPSDEV Check-in Meetings}
Biweekly development team check-in meetings focusing on progress, challenges, concerns, achievements, and overall well-being, with all discussions recorded in the OPSDEV meeting protocols.
\subsection{OPS Review / Consultation Meetings}
Regular OPS team sessions where team members can ask questions, discuss issues, and receive support regarding Anevis products.
\subsection{Quarterly Reviews}
Quarterly review meetings are held where each team presents its current progress, key achievements, and ongoing activities in a structured, sequential manner. The meetings serve to provide an overview of the overall project status across all teams.
\subsection{Internship Feedback Sessions}
Feedback sessions were conducted at regular intervals throughout the internship to exchange mutual feedback, review overall progress, and discuss expectations and satisfaction levels from both sides.

\section{Onboarding}

\subsection{Onboarding Presentation from HR}
The onboarding presentation provided an introduction to Anevis and covered its history, organizational structure, office environment, and company culture. It also included information about meetings, team events, internal communication practices, as well as an overview of the company's products, tools, and working guidelines.

\subsection{OPSDEV Training Exercises}

\begin{itemize}
\item \textit{Exercise Streams:} A set of exercises designed to strengthen proficiency with Java Streams, reflecting their extensive use within the internal system and covering common stream operations and data processing patterns.
\item \textit{Exercise Databases:} Practical exercises focused on gaining experience with MongoDB and PostgreSQL, which are the primary databases used within the internal system, including querying, data modelling, and basic database operations.
\item \textit{Exercise Spring JPA:} An exercise focused on Spring, Hibernate, JPA, SQL, and Java. The task involved generating pie and ring charts with JFreeChart, exporting them to PDF, applying clean OO design principles, and supporting the solution with unit tests, streams, and lambdas.
\item \textit{Exercise Factsheet Project:} A practical factsheet generation exercise simulating a syntetic client project. The task involved working with the internal repositories like document engine, setup scripts; and Git workflows, real databases, JFreeChart reporting components to produce a complete factsheet in a dedicated development branch.

\end{itemize}

\section{OPSDEV - Daily Development Tasks}

\subsection*{OPSDEVBOX Tasks}

OPSDEVBOX taskleri automatisierungta onemli rol oynayan sirketin kendi internal sistemi icin implemente edilen tasklerdir.
Ornegin musteriden gelen datalarin veri bankasina nasil kaydedilecegi, database structurelari update edilmesi yani repositorylerdeki genel duzenlemeler bu task basligi altinda yapilir 


\begin{table}[H]
\centering
\begin{tabular}{p{4cm} p{10cm}}
\toprule
\textbf{Task} & \textbf{Description} \\
\midrule

\texttt{OPSDEVBOX-651} & Add region ultimo mapping to the project setup \\
\texttt{OPSDEVBOX-678} & Make category mandatory in project setup parser \\
\texttt{OPSDEVBOX-679} & Fix marketing classifications with empty identifiers \\
\texttt{OPSDEVBOX-662} & Move series from definitions data to series mappings \\
\texttt{**LLBSWISS-255} & Implement parser to insert valor \\

\bottomrule
\end{tabular}
\caption{OPSDEVBOX Tasks Overview}
\label{tab:opsdevbox-tasks}
\end{table}




\subsection*{BANTLEON Fast Charting Project}

Fast Charting Projesi, Bantleon sirketine ait esg report projesi icin ops team tarafindan hizli ve efektif bir sekilde chart uretimini saglayan databasele baglantisi olmayan pratik yardimci bir tool olarak gelistirilmistir.
Özetle bu class ile chartlar customization ile dilenildigi sekilde visualize edilebilir.

Iste bazi chart ornekleri:

\begin{figure}
    
\end{figure}

\begin{table}[H]
\centering
\begin{tabular}{p{4cm} p{10cm}}
\toprule
\textbf{Task} & \textbf{Description} \\
\midrule
\texttt{BANTLEON-249} & fast-charting-for-bantleon-esg-report \\
\bottomrule
\end{tabular}
\end{table}


\subsection*{Monitoring-Reporter}

Sirkette internal sistem kontrolunu saglamak ve bunlari dilenilen araliklarla mailde raporlastirmak icin bir servis bulunmaktadir ve burada ihtiyaca gore duzenli raporlar uretilebilmektedir.
Burada benim gorevim, sirketin databaseinde bulunan project componenti olmayan projeleri saptamak ve onlarin bilgisini bir tablo halinde mail olarak yonlendirmekti.
Burada internal sistemde zaten hali hazirda var olan bir Class i extend ederek customize edip bir reporter implemente ettim

\begin{table}[H]
\centering
\begin{tabular}{p{4cm} p{10cm}}
\toprule
\textbf{Task} & \textbf{Description} \\
\midrule

\texttt{OPSDEVBOX-680} & monitor-dead-projects \\
\bottomrule
\end{tabular}

\label{tab:opsdevbox-tasks}
\end{table}

\subsection*{Customer-specific Tasks}

Customer specific taskler, pm in musterilerden aldiklari geri bildirim ve taleplere yonelik olusturulan cesitli tasklerdir. 
Bunlara ornek olarak bir musterinin adinin, product adinin degistirilmesi, chartlarin legendinin 2 sutunlu yapilmasi, proje componenti silinmesi, static data update edilmesi, basketlerin ve pricelarin update edilmesi, vb.   
Özetle internship projecte oranla kapsami cok daha kucuk olan yan taskler. 


\section{Additional}
\subsection*{Confluence Article about JFreeCharts}
\subsection*{ISO 27001 Certification Training}
\subsection*{General Data Protection Awareness Training}
\subsection*{Anecon Support}
The perfect opportunity for networking in the financial service industry
Yilda bir kere duzenlenen, cesitli banka ve fon sirketlerinden konusmacilarinin oldugu, \dots
18 Haziran 2026 dan gerceklesen bu konferansta city helper olarak gorev aldim 
\subsection*{Cleaning Service (maybe)}

\section{Given Materials}
The following books were used as reference material during the internship:


\begin{itemize}
    \item Robert C. Martin: \textit{Clean Code: A Handbook of Agile Software Craftsmanship}
    \item Raoul-Gabriel Urma, Mario Fusco, Alan Mycroft: \textit{Java 8 in Action}
\end{itemize}
